% !TEX program = pdflatex
% !TEX encoding = UTF-8 Unicode
% Supervisor meeting, 12 Nov 2026 (DS500). Informal working deck, no title or intro: six slides
% to talk along (status + demo, architecture, validation, server and scaling, our questions,
% next steps) and a short backup. Frozen after the meeting; the next meeting gets a new deck
% (docs/presentations.md, "Freeze rule"). Every number comes from facts.tex next to this file.
% Open items are marked as open; the hosting split is a proposal (PR #64), the Germany outlook
% is a vision of the team (Juju, 12 Nov 2026), not part of the docs' plan.
\documentclass[aspectratio=169,10pt]{beamer}
\usetheme{reach}
% Every number shown in the supervisor deck (12 Nov 2026) comes from here, with the run or
% document it was taken from. Frozen with the deck; never type a number into the deck.
% Rule from docs/notes.md N13: always name profile and budget next to a share.

% ---- validation report (docs/validation_report.md, run 20261106T142630Z) -------------------
\newcommand{\factValRunId}{20261106T142630Z}
\newcommand{\factVOneOvercount}{18\,\%}      % V1: 585 polygon-only pairs = 18.3 % of matrix pairs, walking 3.39
\newcommand{\factVOneBarrier}{47\,\%}        % V1: polygon-only pairs across Neckar/rail/main road
\newcommand{\factVThreeNear}{80--98\,\%}     % V3: picks within 500 m of a removed facility, mode \"created only\" (all deficits: 13--39 %)
\newcommand{\factVThreeRandom}{8--36\,\%}    % V3: same for random picks
\newcommand{\factVFourGap}{2.8\,\%}          % V4: max gap greedy vs. MILP (90 instances)
\newcommand{\factVFourMean}{0.17\,\%}        % V4: mean gap
\newcommand{\factVSixChange}{1.9\,\%}        % V6: residents x categories changing class, walking
\newcommand{\factVEightMoran}{0.72--0.88}    % V8: Moran's I on cells
\newcommand{\factVElevenChange}{3.3--3.6\,\%}% V11: residents x categories changing class

% ---- supervisor meeting 12 Nov 2026 --------------------------------------------------------
% Current dev run: data/runs/20261110T233930Z/validate_results/report_results.md (status ok,
% published, local laptop; 15 categories after PR #60). OSM snapshot from validate/report_inputs.md.
% The seven Core shares equal those of run 20261110T223016Z (interim deck; same OSM snapshot).
\newcommand{\factCurRunId}{20261110T233930Z}
\newcommand{\factCurOsmDate}{8 Nov 2026}
\newcommand{\factCurCells}{1,699}
\newcommand{\factCurResidents}{90,355}
\newcommand{\factCurCategories}{15}
\newcommand{\factCurFacilities}{1,238}            % facilities in the routing area (validate)
\newcommand{\factCurRecs}{224}                     % recommendations_consistent, 3 profiles
\newcommand{\factCurHousing}{830}                  % housing_consistent
\newcommand{\factCurChecks}{15}                    % checks in report part 2
\newcommand{\factCurFalseNeg}{0 of 1,200}          % false_negatives
\newcommand{\factCurMonoCells}{10 of 1,699}        % monotonic_walking_500_ge_walking_339 (fail > 1 % = 16 cells)
\newcommand{\factCurCountRuntime}{254\,s}          % count stage, docs/notes.md N16
% Share of residents per class [%], walking 3.39 km/h, 15 min, run \factCurRunId; sorted by lack.
% ext = extended category (methodology §2, indicative).
\pgfplotstableread{
category                lack  insufficient  sufficient
{Gym (ext.)}            74.2  0.0           25.8
{Senior care (ext.)}    60.6  0.0           39.4
{Library (ext.)}        58.9  0.0           41.1
{Sports centre (ext.)}  49.3  0.0           50.7
{Post office (ext.)}    48.2  0.0           51.8
{Pharmacy}              42.9  0.0           57.1
{Dentist (ext.)}        40.3  0.0           59.7
{Grocery}               32.8  26.2          41.0
{Parks}                 30.8  0.0           69.2
{School}                29.4  0.0           70.6
{Bank \& ATM (ext.)}    23.4  0.0           76.6
{Bakery (ext.)}         21.7  0.0           78.3
{Doctors}               17.5  23.1          59.4
{Childcare}             6.5   14.9          78.6
{Playgrounds}           4.2   0.0           95.8
}\factCurClassesWalk

% Decisions with evidence (docs named per line)
\newcommand{\factUrbanUnscored}{8\,\%}            % architecture.md v0.12 change log, notes N3
\newcommand{\factMatrixAlgoDiff}{2.5\,\%}         % architecture.md v0.15 change log, notes N8
\newcommand{\factPlaygroundsPrivate}{68}          % notes N14: 58 private + 10 customers playgrounds in the city
\newcommand{\factCellsPayload}{59\,KB}            % implementation_plan.md P6: /api/cells gzip, one layer

% Germany outlook: the nationwide Zensus file every run already downloads
% (data/raw/<run>/zensus/Zensus2022_Bevoelkerungszahl_100m-Gitter.csv, counted 12 Nov 2026:
% rows without header, sum of "Einwohner"; values are perturbed by the cell-key method).
\newcommand{\factDeCells}{3,088,037}
\newcommand{\factDeResidents}{82.6\,million}      % 82,570,995
\newcommand{\factDeCellFactor}{1,800}             % 3,088,037 / 1,699 = 1,818
\newcommand{\factDeCountDays}{5}                  % naive: 254 s x 1,818 = 128 h = 5.3 days, one laptop

% Cycling in the recommender, measured 9 Nov 2026 (docs/notes.md N12, W6), macOS/Docker Desktop.
\newcommand{\factCyclingValhallaPeak}{4.0\,GiB}   % + cycling_15, 100 m candidates
\newcommand{\factNoCyclingValhallaPeak}{1.2\,GiB} % walking + wheelchair only
\newcommand{\factCyclingValhallaLimit}{2.5\,GiB}  % limit E6
\newcommand{\factCyclingZones}{12}                % zones cycling would add, for ~130 cells with a deficit
% V2 gaps (docs/notes.md N13)
\newcommand{\factVTwoKitaMissing}{13 of 104}      % Kitas of the city's list without OSM match
\newcommand{\factVTwoSchoolMissing}{2 of 38}      % schools without OSM match

% !TEX root = ../2026-11-12-supervisor.tex
% Schematic isochrone scene, used twice (same streets, river and facilities):
%   \isochronedefinition  what an isochrone is (5/10/15 min, network shape vs. 850 m circle)
%   \isochronetrap        why we do not count with the polygon (methodology §6 V1)
% Scale: 1 unit = 200 m. Walking 3.39 km/h = 56.5 m/min, so 15 min = 850 m = 4.25 units
% along streets. Network times of the facilities (Manhattan paths on the drawn streets):
%   F1 (2.6,4.6): 460 m = 8 min    F2 (6.8,4.6): 780 m = 14 min
%   F3 (6.8,5.5): 960 m = 17 min (inside the 850 m circle, outside the network contour)
%   Fs (4.6,0.6): 1,660 m via the bridge at x = 6.8 = 29 min
% The origin O is at (3.9,3.6). Schematic, not a real place.

\newcommand{\isoscene}{%
  % river (only crossing: the bridge at x = 6.8)
  \fill[river] (-0.2,2.05) .. controls (3,2.1) and (6,1.95) .. (10.2,1.85)
    -- (10.2,1.25) .. controls (6,1.35) and (3,1.5) .. (-0.2,1.45) -- cycle;
  \node[font=\tiny\itshape,text=black!55] at (1.2,1.72) {Neckar};
  % streets north of the river
  \foreach \y in {2.6,3.6,4.6,5.5} \draw[street,line width=1.1pt] (0.2,\y) -- (9.8,\y);
  \foreach \x in {1.2,2.6,3.9,5.3,8.6} \draw[street,line width=1.1pt] (\x,2.35) -- (\x,5.85);
  % street over the bridge and streets south of the river
  \draw[street,line width=1.1pt] (6.8,0.15) -- (6.8,5.85);
  \draw[black!45,line width=2.2pt] (6.8,1.3) -- (6.8,2.05);  % bridge deck
  \draw[street,line width=1.1pt] (0.2,0.5) -- (9.8,0.5);
  \foreach \x in {2.0,4.6,8.6} \draw[street,line width=1.1pt] (\x,0.15) -- (\x,1.1);
}

% Diamond-shaped contour of network distance r (units) around the origin.
\newcommand{\isocontour}[2]{% #1 = radius in units, #2 = tikz options
  \path[#2,rounded corners=4pt] (3.9-#1,3.6) -- (3.9,3.6+#1) -- (3.9+#1,3.6)
    -- (3.9,3.6-#1) -- cycle;}

\newcommand{\isoorigin}{%
  \fill[accent] (3.9,3.6) circle (0.17);
  \draw[white,line width=0.8pt] (3.9,3.6) circle (0.17);}

% \isofacility{x,y}{label}{reached: true/false}{label position: above, right, below left, ...}
\newcommand{\isofacility}[4]{%
  \ifstrequal{#3}{true}%
    {\def\isofill{ink}}{\def\isofill{white}}%
  \node[draw=ink,fill=\isofill,line width=0.8pt,minimum size=2.2mm,inner sep=0pt,
        label={[font=\tiny\bfseries,inner sep=1.5pt,fill=white,fill opacity=0.85,
                text opacity=1]#4:#2}] at (#1) {};}

\newcommand{\isochronedefinition}{%
  \begin{tikzpicture}[scale=0.74]
    \clip (0,0) rectangle (10,6);
    \isoscene
    % true network isochrones stop at the river
    \begin{scope}
      \clip (0,2.2) rectangle (10,6);
      \isocontour{4.25}{fill=sufficient!35,draw=iso,line width=0.9pt}
      \isocontour{2.83}{fill=sufficient!65,draw=iso,line width=0.9pt}
      \isocontour{1.42}{fill=sufficient,draw=iso,line width=0.9pt}
    \end{scope}
    \isoscene  % streets on top of the fills
    \begin{scope}
      \clip (0,2.2) rectangle (10,6);
      \isocontour{4.25}{draw=iso,line width=0.9pt}
      \isocontour{2.83}{draw=iso,line width=0.9pt}
      \isocontour{1.42}{draw=iso,line width=0.9pt}
    \end{scope}
    % straight-line circle of the same 850 m, for comparison
    \draw[muted,dashed,line width=0.7pt] (3.9,3.6) circle (4.25);
    \node[font=\tiny,text=muted,anchor=south west,fill=white,fill opacity=0.85,text opacity=1,
          inner sep=1pt] at (1.25,0.1) {850 m straight line};
    \foreach \r/\t in {1.42/5,2.83/10,4.25/15}
      \node[font=\tiny\bfseries,text=iso,anchor=south,inner sep=1pt] at (3.9+\r-0.25,3.62)
        {\t};
    \isofacility{2.6,4.6}{8 min}{true}{above}
    \isofacility{6.8,4.6}{14 min}{true}{above}
    \isofacility{6.8,5.5}{17 min}{false}{right}
    \isoorigin
    \node[font=\tiny\bfseries,text=accent,anchor=north east,inner sep=2pt] at (3.85,3.5)
      {origin};
  \end{tikzpicture}}

\newcommand{\isochronetrap}{%
  \begin{tikzpicture}[scale=0.74]
    \clip (0,0) rectangle (10,6);
    \isoscene
    % simplified polygon as drawn by an isochrone service: it spills over the river
    \isocontour{4.25}{fill=sufficient!45,fill opacity=0.8}
    \isoscene
    \isocontour{4.25}{draw=iso,line width=0.9pt}
    % the real way to Fs: east to the bridge, across, back west
    \draw[accent,line width=1.4pt,dashed,-{Stealth[length=2.2mm]}]
      (3.9,3.6) -- (3.9,2.6) -- (6.8,2.6) -- (6.8,0.5) -- (4.75,0.5) -- (4.62,0.55);
    \isofacility{6.8,4.6}{14 min}{true}{above}
    \isofacility{2.6,4.6}{8 min}{true}{above}
    \isofacility{4.6,0.6}{29 min}{false}{left}
    \node[font=\tiny\bfseries,text=accent,align=left,anchor=west] at (7.0,1.1)
      {inside the polygon,\\but 29 min on foot};
    \isoorigin
  \end{tikzpicture}}

% !TEX root = ../2026-11-12-supervisor.tex
% System figures (architecture.md §3, §7):
%   \pipelinefigure      stages of one run, from open data to the published `current` views
%   \architecturefigure  components at runtime (README "Architecture in one picture")

\tikzset{
  stage/.style={draw=ink,fill=white,line width=0.6pt,rounded corners=1.5pt,
    font=\scriptsize,minimum height=7mm,minimum width=15mm,align=center,inner sep=1pt},
  gate/.style={stage,draw=accent,line width=1.1pt},
  store/.style={draw=muted,fill=black!5,font=\scriptsize,minimum height=6mm,
    minimum width=16mm,align=center,rounded corners=1.5pt},
  flow/.style={-{Stealth[length=1.6mm]},line width=0.6pt,ink},
  link/.style={line width=0.5pt,muted,dashed},
}

\newcommand{\pipelinefigure}{%
  \begin{tikzpicture}[node distance=2mm]
    \node[stage] (acq) {acquire};
    \node[gate,right=of acq] (val) {validate\\inputs};
    \node[stage,right=of val] (til) {tiles};
    \node[stage,right=of til] (cnt) {count};
    \node[stage,right=of cnt] (sco) {score};
    \node[stage,right=of sco] (rec) {recommend};
    \node[gate,right=of rec] (vre) {validate\\results};
    \node[stage,right=of vre] (pub) {publish};
    \foreach \a/\b in {acq/val,val/til,til/cnt,cnt/sco,sco/rec,rec/vre,vre/pub}
      \draw[flow] (\a) -- (\b);
    % inputs
    \node[store,above=7mm of acq,minimum width=24mm] (src) {OSM $\cdot$ GTFS $\cdot$ Zensus};
    \draw[flow] (src) -- (acq);
    % services
    \node[store,below=7mm of cnt] (vh) {Valhalla};
    \foreach \s in {til,cnt,rec,vre} \draw[link] (\s) -- (vh);
    \node[store,below=7mm of pub] (db) {PostGIS\\\texttt{current}};
    \draw[flow] (pub) -- (db);
    % gate note
    \node[font=\scriptsize,text=accent,align=left,anchor=south west] at ($(val.north)+(0,7mm)$)
      (gn) {a failed check stops the run:\\the app keeps the last good results};
    \draw[accent,line width=0.6pt] (val.north) -- ++(0,6.5mm);
    \draw[accent,line width=0.6pt] (vre.north) |- ($(gn.east)+(1mm,0)$);
  \end{tikzpicture}}

% Components at runtime and the batch path (architecture.md §3, §7, §8): open data -> pipeline
% (one command) -> Valhalla and PostGIS on one server -> FastAPI -> Caddy -> browser.
\newcommand{\architecturefigure}{%
  \begin{tikzpicture}[
      box/.style={draw=ink,fill=white,line width=0.6pt,rounded corners=1.5pt,font=\scriptsize,
        align=center,inner sep=1.5mm,minimum height=10mm},
      svc/.style={box,fill=black!5,draw=muted},
      lab/.style={font=\tiny,text=muted,align=center}]
    \node[svc,text width=19mm] (src) at (0,0)
      {\textbf{Open data}\\OSM (Geofabrik)\\GTFS (NVBW)\\Zensus 2022 grid};
    \node[box,text width=27mm] (pipe) at (30mm,0)
      {\textbf{Pipeline} (\texttt{make refresh})\\acquire $\cdot$ {\color{accent}validate}
      $\cdot$ tiles\\count $\cdot$ score $\cdot$ housing\\recommend $\cdot$
      {\color{accent}validate results}\\publish};
    \node[svc,text width=21mm] (vh) at (63mm,10mm) {\textbf{Valhalla}\\matrix, isochrone,
      route};
    \node[svc,text width=21mm] (db) at (63mm,-10mm) {\textbf{PostGIS}\\one schema per run,
      view \texttt{current}};
    \node[box,text width=17mm] (api) at (91mm,0) {\textbf{FastAPI}\\read + live
      endpoints};
    \node[box,text width=17mm] (caddy) at (114mm,0) {\textbf{Caddy}\\HTTPS, static web
      build};
    \node[box,text width=17mm] (web) at (114mm,19mm) {\textbf{Browser}\\React +
      MapLibre};
    \begin{scope}[on background layer]
      \node[draw=muted,dashed,rounded corners=2pt,inner sep=2.5mm,
        fit=(pipe)(vh)(db)(api)(caddy)] (srv) {};
    \end{scope}
    \node[lab,anchor=south west] at (srv.north west) {one server, Docker Compose (host:
      open question 1)};
    \draw[flow] (src) -- (pipe);
    \draw[flow] (pipe) -- (vh);
    \draw[flow] (pipe) -- (db);
    \draw[flow] (api) -- (vh);
    \draw[flow] (api) -- (db);
    \draw[flow] (caddy) -- node[lab,above]{\texttt{/api}} (api);
    \draw[flow] (web) -- node[lab,right]{HTTPS} (caddy);
  \end{tikzpicture}}

% !TEX root = ../2026-11-12-supervisor.tex
\usetikzlibrary{fit,backgrounds}
% Hosting figures for the supervisor discussion (docs/architecture.md §3, §11 = current design;
% docs/deployment_proposal.md = proposal of PR #64, not approved).
%   \hostingcurrent   binding design: one small server runs Caddy, API, PostGIS, Valhalla, pipeline
%   \hostingproposed  proposal: static frontend on Cloudflare Pages, backend on one host behind a
%                     Cloudflare Tunnel, team access over a private network (candidate: Tailscale)
% Both are drawn about 66 mm wide so that they fit next to each other in two columns.

\tikzset{
  hnode/.style={draw=ink,line width=0.6pt,align=center,font=\scriptsize,inner sep=1.6mm,
    minimum height=6.5mm},
  hprop/.style={hnode,dashed,draw=accent},
  hhost/.style={draw=ink,line width=0.6pt,rounded corners=1pt,inner sep=2mm},
  harrow/.style={-{Stealth[length=1.8mm]},line width=0.6pt,draw=ink},
  hlabel/.style={font=\tiny,text=muted},
}

\newcommand{\hostingcurrent}{%
  \begin{tikzpicture}
    \node[hnode] (v) at (0,0) {Visitors};
    \node[hnode,minimum width=30mm] (c) at (0,-1.3) {Caddy: web build, HTTPS, \texttt{/api}};
    \node[hnode] (api) at (-1.9,-2.7) {FastAPI};
    \node[hnode] (db) at (0,-2.7) {PostGIS};
    \node[hnode] (va) at (1.9,-2.7) {Valhalla};
    \node[hnode,minimum width=30mm] (pl) at (0,-3.75) {Pipeline, tile build (\texttt{make refresh})};
    \begin{pgfonlayer}{background}
      \node[hhost,fit=(c)(api)(db)(va)(pl),label={[font=\tiny,text=muted,anchor=south east,inner sep=0.6mm]
        north east:one small server}] {};
    \end{pgfonlayer}
    \draw[harrow] (v) -- (c);
    \draw[harrow] (c) -- (api);
  \end{tikzpicture}}

\newcommand{\hostingproposed}{%
  \begin{tikzpicture}
    \node[hnode] (v) at (0,0) {Visitors};
    \node[hprop,minimum width=24mm] (p) at (-1.7,-1.1) {Cloudflare Pages\\static web build};
    \node[hprop,minimum width=24mm] (t) at (1.7,-1.1) {Cloudflare Tunnel\\API hostname};
    \node[hnode] (api) at (-1.9,-2.6) {FastAPI};
    \node[hnode] (db) at (0,-2.6) {PostGIS};
    \node[hnode] (va) at (1.9,-2.6) {Valhalla};
    \node[hnode,minimum width=30mm] (pl) at (0,-3.55) {Pipeline, tile build};
    \begin{pgfonlayer}{background}
      \node[hhost,fit=(api)(db)(va)(pl),label={[font=\tiny,text=muted,anchor=south east,inner sep=0.6mm]
        north east:one backend host (ProDesk or VPS)}] (host) {};
    \end{pgfonlayer}
    \node[hprop,minimum width=22mm] (ts) at (3.0,-4.45) {Team: private\\network (Tailscale?)};
    \draw[harrow] (v) -- (p);
    \draw[harrow] (v) -- (t);
    \draw[harrow] (t) -- (api);
    \draw[harrow,dashed,draw=accent] (ts) -- (pl);
  \end{tikzpicture}}


% Ragged-right table column (no stretched word spaces in narrow cells).
\newcolumntype{P}[1]{>{\raggedright\arraybackslash}p{#1}}

\title[reach $\cdot$ Supervisor meeting]{reach: where we stand}
\author{Rico $\cdot$ Socke $\cdot$ Juju}
\institute{University of Tübingen, winter term 2026/27}
\date[12 Nov 2026]{12 November 2026}

\begin{document}

% ============================================================================================
% Speaker: TBD
\begin{frame}{reach already runs end to end for Tübingen; we show you the app live}
  \begin{columns}[T,onlytextwidth]
    \column{0.30\textwidth}
      \keyfigure{\factCurCells}{populated 100\,m cells (\factCurResidents{} residents)}
      \par\bigskip
      \keyfigure{\factCurCategories{} $\times$ 4}{categories $\times$ profiles (walking at
        two speeds, cycling, wheelchair)}
      \par\bigskip
      \keyfigure{9 of 11}{validation experiments done}
    \column{0.66\textwidth}
      \textbf{What we show in the app}
      \begin{enumerate}\footnotesize
        \item \textbf{Map:} who lacks access to what, per category and profile
        \item \textbf{Click anywhere:} what is reachable in 15 minutes, with times and routes,
              and the nearest missing facility
        \item \textbf{Combined score} 0--100 over the categories you pick
        \item \textbf{Recommendations:} where a new facility would help the most residents
              (\factCurRecs{} zones)
        \item \textbf{Housing potential:} empty cells with good access for new homes
              (\factCurHousing{} candidates)
        \item \textbf{Data page:} sources, snapshot dates, validation report
      \end{enumerate}
      \medskip
      {\footnotesize\color{muted}Runs locally today; pipeline, API and web app are on
      \texttt{dev}.}
  \end{columns}
  \source{run \factCurRunId{} (OSM \factCurOsmDate); validation report (run \factValRunId);
    docs/features.md §1--§5.}
\end{frame}

% ============================================================================================
% Speaker: TBD
\begin{frame}{How it works: one command turns open data into the map}
  \centering
  \architecturefigure\par\medskip
  \raggedright
  \begin{columns}[T,onlytextwidth]
    \column{0.32\textwidth}
      \footnotesize\textbf{Open data only:} OSM, timetables, census grid; one snapshot per
      run.
    \column{0.32\textwidth}
      \footnotesize\textbf{Real travel times:} routing on the street network, not straight
      lines.
    \column{0.32\textwidth}
      \footnotesize\textbf{Checked:} \factCurChecks{} automatic checks; a failure keeps the
      last good run.
  \end{columns}
  \source{docs/architecture.md §2, §3, §7, §8; run \factCurRunId{} (validate\_results report);
    weekly refresh and production server are planned (§11).}
\end{frame}

% ============================================================================================
% Speaker: TBD
\begin{frame}{We tested our method against the obvious alternatives, and it holds up}
  \begin{columns}[T,onlytextwidth]
    \column{0.24\textwidth}
      \keyfigure{$+$\factVOneOvercount}{wrong pairs if we counted with isochrone polygons}
      \par\smallskip
      {\footnotesize Mostly across the Neckar, rail lines and main roads. We count with real
      travel times instead.}
    \column{0.24\textwidth}
      \keyfigure{$\le$\,\factVFourGap}{below the optimum: our fast recommender vs.\ an exact
        solver}
      \par\smallskip
      {\footnotesize 90 test instances; good enough to run in every refresh.}
    \column{0.24\textwidth}
      \keyfigure{\factVTwoKitaRecall}{of the city's Kitas found in OSM}
      \par\smallskip
      {\footnotesize Schools: \factVTwoSchoolRecall. The gaps are known; we can map them in
      OSM.}
    \column{0.24\textwidth}
      \keyfigure{\factVSixChange}{change class if we used buildings instead of cell centres}
      \par\smallskip
      {\footnotesize Walking; the 100\,m grid is precise enough.}
  \end{columns}
  \vfill
  {\footnotesize\color{muted}Still to do: a check against local knowledge and a full
  reproducibility run. All results: backup.}
  \source{validation report V1, V2, V4, V6 (run \factValRunId); walking 3.39\,km/h, 15 min.}
\end{frame}

% ============================================================================================
% Speaker: TBD
\begin{frame}{Tübingen fits on one small server; all of Germany would need a bigger setup}
  \begin{columns}[T,onlytextwidth]
    \column{0.48\textwidth}
      \textbf{Server for Tübingen}
      \begin{itemize}\footnotesize
        \item Counting all \factCurCategories{} categories $\times$ 4 profiles takes
              \factCurCountRuntime{} on a laptop
        \item The routing engine needs up to \factCyclingValhallaPeak{} of memory, so we plan
              for 8\,GB RAM
        \item \textbf{Option A:} rented VPS ($\ge$\,8\,GB); budget 20\,EUR in total, so online
              from about 1 Dec until grading (price still to check)
        \item \textbf{Option B:} our own HP ProDesk (i5, 16\,GB): no cost, but home network and
              uptime
        \item Optional: web frontend separately on Cloudflare Pages (proposal)
      \end{itemize}
    \column{0.48\textwidth}
      \textbf{All of Germany (idea, not planned)}
      \begin{itemize}\footnotesize
        \item The census file we already use covers \factDeCells{} populated cells
              (\factDeResidents{} residents), about \factDeCellFactor{}$\times$ Tübingen
        \item Same grid, same pipeline, one run per municipality
        \item Heavy routing builds on bwUniCluster instead of our server
        \item Hard parts: OSM quality differs by region; villages need other thresholds than
              a city
      \end{itemize}
  \end{columns}
  \source{docs/architecture.md §11, §11.1; docs/deployment\_proposal.md (PR \#64, not approved);
    notes N12, N16 (laptop, run \factCurRunId); Zensus file of run \factCurRunId{} (grid sum).}
\end{frame}

% ============================================================================================
% Speaker: TBD
\begin{frame}{What we would like to discuss with you today}
  \footnotesize\renewcommand{\arraystretch}{1.35}
  \begin{tabular}{@{}P{80mm}P{55mm}@{}}
    \toprule
    Question & Our suggestion \\
    \midrule
    \textbf{Server:} Is our own ProDesk fine, or do you expect a rented VPS? How long does
      the app need to stay online? & Measure one full refresh on the ProDesk, then decide \\
    \textbf{Access:} How would you like to look at the data: app only, or also the database
      (e.g.\ via SSH tunnel)? & Read-only database account \\
    \textbf{Scaling:} Is Tübingen enough, or would a second city or a larger region be
      interesting? May we use bwUniCluster? & Finish Tübingen first; Germany as outlook \\
    \textbf{Method:} Moran's I shows strong clusters on cells (\factVEightMoran). Is that
      enough as the ``adequate method'', or do you expect more? & Optimisation plus spatial
      statistics \\
    \textbf{Limits:} Is it okay to name the known limits openly (wheelchair on slopes, OSM
      gaps, new categories not validated yet)? & Yes, on the data page and in the report \\
    \textbf{Interim on 25 Nov:} What do you expect: live demo, results, or both? & Live demo
      with backup screenshots \\
    \bottomrule
  \end{tabular}
  \source{docs/deployment\_proposal.md (PR \#64); methodology.md §7 (M5); validation report V8
    (run \factValRunId); notes N11, N13, N16; project\_management\_decisions.md §8.}
\end{frame}

% ============================================================================================
% Speaker: TBD
\begin{frame}{Next steps: server live in December, submission on 22 January}
  \begin{columns}[T,onlytextwidth]
    \column{0.48\textwidth}
      \textbf{Plan until the submission}
      \begin{itemize}\footnotesize
        \item \textbf{25 Nov:} interim presentation (\texttt{v0.1.0})
        \item \textbf{from $\sim$1 Dec:} server live, weekly automatic refresh
        \item \textbf{Dec--Jan:} remaining validation (local knowledge, reproducibility);
              final categories and public-transport classes
        \item \textbf{10--15 Jan:} feature freeze
        \item \textbf{22 Jan:} submission (\texttt{v1.0.0}); \textbf{29 Jan:} final
              presentation
      \end{itemize}
    \column{0.48\textwidth}
      \textbf{What we need from today}
      \begin{itemize}\footnotesize
        \item A decision on the server, so we can deploy from 1 Dec
        \item Your view on scope: Tübingen only, or a look beyond
        \item Feedback on the app: what is missing, what is too much?
      \end{itemize}
  \end{columns}
  \source{docs/features.md §9; architecture.md §11; methodology.md M1, M2.}
\end{frame}

% ============================================================================================
\backupslides
\section*{Backup}

\begin{frame}{Backup: validation experiments V1--V11}
  \scriptsize
  \begin{tabular}{@{}llP{95mm}@{}}
    \toprule
    \# & Experiment & Key result \\
    \midrule
    V1 & Matrix vs.\ isochrone polygons & Polygons add 18--20\,\% false pairs, 47\,\% across barriers \\
    V2 & OSM completeness & Childcare recall 88\,\%, precision 98\,\%; schools 95\,\% / 95\,\% \\
    V3 & Recommender back-test & 80--98\,\% of picks within 500\,m of a removed facility (only the removed deficit counts) \\
    V4 & Greedy vs.\ exact MILP & At most 2.8\,\% below the optimum (mean 0.17\,\%) \\
    V5 & Sensitivity ($k$, speed, budget) & Shares move strongly; district rankings stable
      (median Spearman 0.78--0.95) \\
    V6 & Centroid vs.\ building origin & 1.9\,\% (walking) / 2.8\,\% (wheelchair) change class \\
    V7 & Local knowledge & Not yet run \\
    V8 & Spatial statistics & Moran's I 0.72--0.88 on cells; none between districts \\
    V9 & Reproducibility & Not yet run \\
    V10 & Matrix vs.\ single routes & Median $|\Delta|$ 0.5\,s; 0 of 900 false negatives \\
    V11 & With vs.\ without slopes & 3.3--3.6\,\% change class, mostly on hillsides \\
    \bottomrule
  \end{tabular}
  \source{docs/validation\_report.md, run \factValRunId{} (Core categories of that run, parks
    and playgrounds still one category); the extended categories are not validated yet.}
\end{frame}

\begin{frame}{Backup: counting with isochrone polygons would add \factVOneOvercount{}
  unreachable pairs}
  \begin{columns}[T,onlytextwidth]
    \column{0.55\textwidth}
      \centering\isochronetrap
    \column{0.42\textwidth}
      \footnotesize
      \textbf{How we count:} per block of cells one travel-time matrix to all facilities in
      reach; a facility counts if its time is $\le 15$ min. Isochrones are only drawn.
      \par\medskip
      \textbf{Test} (150 cells, walking 3.39\,km/h): polygons add \factVOneOvercount{}
      unreachable pairs, \factVOneBarrier{} of them across the Neckar, rail lines or main
      roads.\par\medskip
      \textbf{Nothing lost:} \factCurFalseNeg{} sampled pairs outside the counts are
      reachable by a single route.
  \end{columns}
  \source{docs/architecture.md §6.3, §6.5; validation report V1 (run \factValRunId); run
    \factCurRunId{} (\texttt{false\_negatives}).}
\end{frame}

\begin{frame}{Backup: hosting, current design vs.\ the Cloudflare proposal}
  \begin{columns}[T,onlytextwidth]
    \column{0.5\textwidth}
      \centering{\footnotesize\bfseries Current design}\par\smallskip
      \hostingcurrent
    \column{0.5\textwidth}
      \centering{\footnotesize\bfseries Proposal (not approved)}\par\smallskip
      \hostingproposed
  \end{columns}
  \source{docs/architecture.md §3, §11 (current); docs/deployment\_proposal.md (proposal, PR
    \#64): backend on the HP ProDesk (i5-8500T, 16\,GB) or a rented VPS ($\ge 8$\,GB).}
\end{frame}

\begin{frame}{Backup: known limits and how we handle them}
  \footnotesize
  \begin{tabular}{@{}P{32mm}P{62mm}P{34mm}@{}}
    \toprule
    Limit & Effect & Proposal \\
    \midrule
    Wheelchair on slopes & Steep sections are not avoided; hillside counts are optimistic
      & Accept, state on data page \\
    OSM gaps & \factVTwoKitaMissing{} Kitas and \factVTwoSchoolMissing{} schools have no OSM
      match & Map them in OSM \\
    New categories & No completeness check yet; consistency check at \factCurMonoCells{}
      cells, fails above 1\,\% & Mark as indicative \\
    Cycling recommender off & Would add only \factCyclingZones{} zones but needs
      \factCyclingValhallaPeak{} of routing memory & Accept, document \\
    Housing candidates & Land use not fully excluded; buildable not known & Decide after
      map check \\
    \bottomrule
  \end{tabular}
  \par\medskip
  {\footnotesize\color{muted}Shares depend strongly on speed and budget: we always name the
  profile and budget next to a share.}
  \source{notes N11, N12, N13, N15, N16; validation report V2, V5 (run \factValRunId); run
    \factCurRunId.}
\end{frame}

\begin{frame}{Backup: we use the official census 100\,m grid instead of H3 hexagons}
  \footnotesize
  \begin{tabular}{@{}P{20mm}P{57mm}P{55mm}@{}}
    \toprule
    & \textbf{Census 100\,m grid (chosen)} & H3 hexagons \\
    \midrule
    Residents & Exact counts as published & Re-gridding spreads residents by area; errors at
      every edge \\
    Identity & Official cell id, joins other census tables (e.g.\ age) & Own index; every join
      needs re-gridding \\
    Scale & Same grid for all of Germany: \factDeCells{} populated cells in the file we already
      download & Global, but residents must be disaggregated first \\
    Trade-off & Squares: 8 neighbours at two distances & Hexagons: 6 equal neighbours, nicer
      for spatial statistics \\
    \bottomrule
  \end{tabular}
  \source{docs/architecture.md §2, §5.2, §12; validation report V6 (run \factValRunId); Zensus
    file of run \factCurRunId.}
\end{frame}

\end{document}
